\section{Method}\label{sec:method}
Define the dataset $\mathcal{D}=\{(x_i,y_i)\}_{i=1}^{N}$ and predictive
model $f_{\theta}$. A generic example training objective is
\begin{equation}\label{eq:objective}
\mathcal{L}(\theta)=\frac{1}{N}\sum_{i=1}^{N}
\ell\bigl(f_{\theta}(x_i),y_i\bigr).
\end{equation}
Explain each variable, modeling choice, and evaluation assumption. Use
\cref{eq:objective} to connect the mathematical formulation to your method.

\subsection{Method overview}
Provide a diagram only when it clarifies the procedure. \Cref{fig:overview}
is an illustrative example.
\begin{figure}[t]
\centering
\begin{tikzpicture}[>=Stealth,node distance=0.32cm,
 every node/.style={draw=uiucline,fill=uiucpale,rounded corners=2pt,
 minimum width=0.86\columnwidth,minimum height=0.58cm,
 align=center,font=\small}]
\node (a) {Observed samples};
\node[below=0.31cm of a] (b) {Prediction model};
\node[below=0.31cm of b] (c) {Held-out evaluation};
\draw[->,line width=0.75pt,draw=illinoisnavy] (a)--(b);
\draw[->,line width=0.75pt,draw=illinoisnavy] (b)--(c);
\end{tikzpicture}
\caption{Illustrative pipeline for demonstrating figure styling.}
\label{fig:overview}
\end{figure}

\subsection{Implementation details}
Report model configurations, train/test splits, hyperparameters, hardware,
random seeds, and code availability. Define baselines consistently.
